%!TeX root=../pridetop.tex
\chapter[Chapter \thechapter]{}
\begin{figure}[t!]
\centering
\includegraphics[width=\linewidth]{6top}
\captionlistentry{Headpiece to Chapter \thechapter}
\end{figure}

\lettrine[lines=6,image=true]{initials/chap6t}{he}  ladies of Longbourn soon waited on those of Netherfield. The visit was returned in due form. Miss Bennet's pleasing manners grew on the good-will of Mrs Hurst and Miss Bingley; and though the mother was found to be intolerable, and the younger sisters not worth speaking to, a wish of being better acquainted with \textit{them} was expressed towards the two eldest. By Jane this attention was received with the greatest pleasure; but Elizabeth still saw superciliousness in their treatment of everybody, hardly excepting even her sister, and could not like them; though their kindness to Jane, such as it was, had a value, as arising, in all probability, from the influence of their brother's admiration. It was generally evident, whenever they met, that he \textit{did} admire her; and to \textit{her} it was equally evident that Jane was yielding to the preference which she had begun to entertain for him from the first, and was in a way to be very much in love; but she considered with pleasure that it was not likely to be discovered by the world in general, since Jane united with great strength of feeling, a composure of temper and an uniform cheerfulness of manner, which would guard her from the suspicions of the impertinent. She mentioned this to her friend, Miss Lucas.

»It may, perhaps, be pleasant,« replied Charlotte, »to be able to impose on the public in such a case; but it is sometimes a disadvantage to be so very guarded. If a woman conceals her affection with the same skill from the object of it, she may lose the opportunity of fixing him; and it will then be but poor consolation to believe the world equally in the dark. There is so much of gratitude or vanity in almost every attachment, that it is not safe to leave any to itself. We can all \textit{begin} freely—a slight preference is natural enough; but there are very few of us who have heart enough to be really in love without encouragement. In nine cases out of ten, a woman had better show \textit{more} affection than she feels. Bingley likes your sister undoubtedly; but he may never do more than like her, if she does not help him on.«

»But she does help him on, as much as her nature will allow. If \textit{I} can perceive her regard for him, he must be a simpleton indeed not to discover it too.«

»Remember, Eliza, that he does not know Jane's disposition as you do.«

»But if a woman is partial to a man, and does not endeavour to conceal it, he must find it out.«

»Perhaps he must, if he sees enough of her. But though Bingley and Jane meet tolerably often, it is never for many hours together; and as they always see each other in large mixed parties, it is impossible that every moment should be employed in conversing together. Jane should therefore make the most of every half hour in which she can command his attention. When she is secure of him, there will be leisure for falling in love as much as she chooses.«

»Your plan is a good one,« replied Elizabeth, »where nothing is in question but the desire of being well married; and if I were determined to get a rich husband, or any husband, I dare say I should adopt it. But these are not Jane's feelings; she is not acting by design. As yet she cannot even be certain of the degree of her own regard, nor of its reasonableness. She has known him only a fortnight. She danced four dances with him at Meryton; she saw him one morning at his own house, and has since dined in company with him four times. This is not quite enough to make her understand his character.«

»Not as you represent it. Had she merely \textit{dined} with him, she might only have discovered whether he had a good appetite; but you must remember that four evenings have been also spent to\-geth\-er—and four evenings may do a great deal.«

»Yes: these four evenings have enabled them to ascertain that they both like Vingt-un better than Commerce, but with respect to any other leading characteristic, I do not imagine that much has been unfolded.«

»Well,« said Charlotte, »I wish Jane success with all my heart; and if she were married to him to-morrow, I should think she had as good a chance of happiness as if she were to be studying his character for a twelvemonth. Happiness in marriage is entirely a matter of chance. If the dispositions of the parties are ever so well known to each other, or ever so similar beforehand, it does not advance their felicity in the least. They always continue to grow sufficiently unlike afterwards to have their share of vexation; and it is better to know as little as possible of the defects of the person with whom you are to pass your life.«

»You make me laugh, Charlotte; but it is not sound. You know it is not sound, and that you would never act in this way yourself.«

Occupied in observing Mr Bingley's attention to her sister, Elizabeth was far from suspecting that she was herself becoming an object of some interest in the eyes of his friend. Mr Darcy had at first scarcely allowed her to be pretty: he had looked at her without admiration at the ball; and when they next met, he looked at her only to criticise. But no sooner had he made it clear to himself and his friends that she had hardly a good feature in her face, than he began to find it was rendered uncommonly intelligent by the beautiful expression of her dark eyes. To this discovery succeeded some others equally mortifying. Though he had detected with a critical eye more than one failure of perfect symmetry in her form, he was forced to acknowledge her figure to be light and pleasing; and in spite of his asserting that her manners were not those of the fashionable world, he was caught by their easy playfulness. Of this she was perfectly unaware: to her he was only the man who made himself agreeable nowhere, and who had not thought her handsome enough to dance with.

He began to wish to know more of her; and, as a step towards conversing with her himself, attended to her conversation with others. His doing so drew her notice. It was at Sir William Lucas's, where a large party were assembled.

»What does Mr Darcy mean,« said she to Charlotte, »by listening to my conversation with Colonel Forster?«

»That is a question which Mr Darcy only can answer.«

»But if he does it any more, I shall certainly let him know that I see what he is about. He has a very satirical eye, and if I do not begin by being impertinent myself, I shall soon grow afraid of him.«


On his approaching them soon afterwards, though without seeming to have any intention of speaking, Miss Lucas defied her friend to mention such a subject to him, which immediately provoking Elizabeth to do it, she turned to him and said,—

»Did not you think, Mr Darcy, that I expressed myself uncommonly well just now, when I was teasing Colonel Forster to give us a ball at Meryton?«

»With great energy; but it is a subject which always makes a lady energetic.«

»You are severe on us.«

»It will be \textit{her} turn soon to be teased,« said Miss Lucas. »I am going to open the instrument, Eliza, and you know what follows.«

»You are a very strange creature by way of a friend!—always wanting me to play and sing before anybody and everybody! If my vanity had taken a musical turn, you would have been invaluable; but as it is, I would really rather not sit down before those who must be in the habit of hearing the very best performers.« On Miss Lucas's persevering, however, she added, »Very well; if it must be so, it must.« And gravely glancing at Mr Darcy, “There is a very fine old saying, which everybody here is of course familiar with—»Keep your breath to cool your porridge,«—and I shall keep mine to swell my song.”

Her performance was pleasing, though by no means capital. After a song or two, and before she could reply to the entreaties of several that she would sing again, she was eagerly succeeded at the instrument by her sister Mary, who having, in consequence of being the only plain one in the family, worked hard for knowledge and accomplishments, was always impatient for display.

\begin{figure}[tbh]
\centering
\includegraphics[width=\linewidth]{6entreaties}
\captionlistentry{The entreaties of several}
\end{figure}

Mary had neither genius nor taste; and though vanity had given her application, it had given her likewise a pedantic air and conceited manner, which would have injured a higher degree of excellence than she had reached. Elizabeth, easy and unaffected, had been listened to with much more pleasure, though not playing half so well; and Mary, at the end of a long concerto, was glad to purchase praise and gratitude by Scotch and Irish airs, at the request of her younger sisters, who with some of the Lucases, and two or three officers, joined eagerly in dancing at one end of the room.

Mr Darcy stood near them in silent indignation at such a mode of passing the evening, to the exclusion of all conversation, and was too much engrossed by his own thoughts to perceive that Sir William Lucas was his neighbour, till Sir William thus began:—

»What a charming amusement for young people this is, Mr Darcy! There is nothing like dancing, after all. I consider it as one of the first refinements of polished societies.«

»Certainly, sir; and it has the advantage also of being in vogue amongst the less polished societies of the world: every savage can dance.«

Sir William only smiled. »Your friend performs delightfully,« he continued, after a pause, on seeing Bingley join the group; »and I doubt not that you are an adept in the science yourself, Mr Darcy.«

»You saw me dance at Meryton, I believe, sir.«

»Yes, indeed, and received no inconsiderable pleasure from the sight. Do you often dance at St James's?«

»Never, sir.«

»Do you not think it would be a proper compliment to the place?«

»It is a compliment which I never pay to any place if I can avoid it.«

»You have a house in town, I conclude?«

Mr Darcy bowed.

»I had once some thoughts of fixing in town myself, for I am fond of superior society; but I did not feel quite certain that the air of London would agree with Lady Lucas.«

He paused in hopes of an answer: but his companion was not disposed to make any; and Elizabeth at that instant moving towards them, he was struck with the notion of doing a very gallant thing, and called out to her,—

»My dear Miss Eliza, why are not you dancing? Mr Darcy, you must allow me to present this young lady to you as a very desirable partner. You cannot refuse to dance, I am sure, when so much beauty is before you.« And, taking her hand, he would have given it to Mr Darcy, who, though extremely surprised, was not unwilling to receive it, when she instantly drew back, and said with some discomposure to Sir William,—

»Indeed, sir, I have not the least intention of dancing. I entreat you not to suppose that I moved this way in order to beg for a partner.«

Mr Darcy, with grave propriety, requested to be allowed the honour of her hand, but in vain. Elizabeth was determined; nor did Sir William at all shake her purpose by his attempt at persuasion.

»You excel so much in the dance, Miss Eliza, that it is cruel to deny me the happiness of seeing you; and though this gentleman dislikes the amusement in general, he can have no objection, I am sure, to oblige us for one half hour.«

»Mr Darcy is all politeness,« said Elizabeth, smiling.

»He is, indeed: but considering the inducement, my dear Miss Eliza, we cannot wonder at his complaisance; for who would object to such a partner?«

Elizabeth looked archly, and turned away. Her resistance had not injured her with the gentleman, and he was thinking of her with some complacency, when thus accosted by Miss Bingley,—

»I can guess the subject of your reverie.«

»I should imagine not.«

»You are considering how insupportable it would be to pass many evenings in this manner,—in such society; and, indeed, I am quite of your opinion. I was never more annoyed! The insipidity, and yet the noise—the nothingness, and yet the self-importance, of all these people! What would I give to hear your strictures on them!«

»Your conjecture is totally wrong, I assure you. My mind was more agreeably engaged. I have been meditating on the very great pleasure which a pair of fine eyes in the face of a pretty woman can bestow.«

Miss Bingley immediately fixed her eyes on his face, and desired he would tell her what lady had the credit of inspiring such reflections. Mr Darcy replied, with great intrepidity,—

»Miss Elizabeth Bennet.«

»Miss Elizabeth Bennet!« repeated Miss Bingley. »I am all astonishment. How long has she been such a favourite? and pray when am I to wish you joy?«

»That is exactly the question which I expected you to ask. A lady's imagination is very rapid; it jumps from admiration to love, from love to matrimony, in a moment. I knew you would be wishing me joy.«

»Nay, if you are so serious about it, I shall consider the matter as absolutely settled. You will have a charming mother-in-law, indeed, and of course she will be always at Pemberley with you.«

He listened to her with perfect indifference, while she chose to entertain herself in this manner; and as his composure convinced her that all was safe, her wit flowed along.